\documentclass[10pt,a4paper]{amsart}
\usepackage{amsmath,amssymb,amsthm,graphicx}
% For commutative diagrams we use Xy-pic
\usepackage[all]{xy}

% We use 2cell for 2-commutative diagrams.
\xyoption{2cell}
% Theorem environments.
%
\theoremstyle{plain}
\newtheorem{theorem}[subsection]{Theorem}
\newtheorem{proposition}[subsection]{Proposition}
\newtheorem{lemma}[subsection]{Lemma}

\theoremstyle{definition}
\newtheorem{definition}[subsection]{Definition}
\newtheorem{example}[subsection]{Example}
\newtheorem{exercise}[subsection]{Exercise}
\newtheorem{situation}[subsection]{Situation}

\theoremstyle{remark}
\newtheorem{remark}[subsection]{Remark}
\newtheorem{remarks}[subsection]{Remarks}

\numberwithin{equation}{subsection}

% Macros
%
\def\lim{\mathop{\mathrm{lim}}\nolimits}
\def\colim{\mathop{\mathrm{colim}}\nolimits}
\def\Spec{\mathop{\mathrm{Spec}}}
\def\Hom{\mathop{\mathrm{Hom}}\nolimits}
\def\Ext{\mathop{\mathrm{Ext}}\nolimits}
\def\SheafHom{\mathop{\mathcal{H}\!\mathit{om}}\nolimits}
\def\SheafExt{\mathop{\mathcal{E}\!\mathit{xt}}\nolimits}
\def\Sch{\mathit{Sch}}
\def\Mor{\mathop{\mathrm{Mor}}\nolimits}
\def\Ob{\mathop{\mathrm{Ob}}\nolimits}
\def\Sh{\mathop{\mathit{Sh}}\nolimits}
\def\NL{\mathop{N\!L}\nolimits}
\def\CH{\mathop{\mathrm{CH}}\nolimits}
\def\proetale{{pro\text{-}\acute{e}tale}}
\def\etale{{\acute{e}tale}}
\def\QCoh{\mathit{QCoh}}
\def\Ker{\mathop{\mathrm{Ker}}}
\def\Im{\mathop{\mathrm{Im}}}
\def\Coker{\mathop{\mathrm{Coker}}}
\def\Coim{\mathop{\mathrm{Coim}}}

% Boxtimes
%
\DeclareMathSymbol{\boxtimes}{\mathbin}{AMSa}{"02}

%
% Macros for moduli stacks/spaces
%
\def\QCohstack{\mathcal{QC}\!\mathit{oh}}
\def\Cohstack{\mathcal{C}\!\mathit{oh}}
\def\Spacesstack{\mathcal{S}\!\mathit{paces}}
\def\Quotfunctor{\mathrm{Quot}}
\def\Hilbfunctor{\mathrm{Hilb}}
\def\Curvesstack{\mathcal{C}\!\mathit{urves}}
\def\Polarizedstack{\mathcal{P}\!\mathit{olarized}}
\def\Complexesstack{\mathcal{C}\!\mathit{omplexes}}
% \Pic is the operator that assigns to X its picard group, usage \Pic(X)
% \Picardstack_{X/B} denotes the Picard stack of X over B
% \Picardfunctor_{X/B} denotes the Picard functor of X over B
\def\Pic{\mathop{\mathrm{Pic}}\nolimits}
\def\Picardstack{\mathcal{P}\!\mathit{ic}}
\def\Picardfunctor{\mathrm{Pic}}
\def\Deformationcategory{\mathcal{D}\!\mathit{ef}}
\usepackage{cleveref}
\crefname{originallemma}{original lemma}{original lemmas}
\Crefname{originallemma}{Original lemma}{Original lemmas}
\makeatletter
\newlabel{algebra-lemma-finite-extension-dim-1}{{02MJ}{source}{}{sourcepointer-1}{}}
\newlabel{algebra-lemma-finite-extension-dim-1@cref}{{[originallemma][0][]02MJ}{[1][0][]source}{}{}{}}
\makeatother
\crefname{originallemma}{original lemma}{original lemmas}
\Crefname{originallemma}{Original lemma}{Original lemmas}
\makeatletter
\newlabel{lemma-compare-periodic-lengths}{{0EA9}{source}{}{sourcepointer-2}{}}
\newlabel{lemma-compare-periodic-lengths@cref}{{[originallemma][0][]0EA9}{[1][0][]source}{}{}{}}
\makeatother
\crefname{originallemma}{original lemma}{original lemmas}
\Crefname{originallemma}{Original lemma}{Original lemmas}
\makeatletter
\newlabel{lemma-length-multiplication}{{02QF}{source}{}{sourcepointer-3}{}}
\newlabel{lemma-length-multiplication@cref}{{[originallemma][0][]02QF}{[1][0][]source}{}{}{}}
\makeatother
\crefname{originallemma}{original lemma}{original lemmas}
\Crefname{originallemma}{Original lemma}{Original lemmas}
\makeatletter
\newlabel{lemma-multiply-period-length}{{0EAC}{source}{}{sourcepointer-4}{}}
\newlabel{lemma-multiply-period-length@cref}{{[originallemma][0][]0EAC}{[1][0][]source}{}{}{}}
\makeatother
\crefname{originallemma}{original lemma}{original lemmas}
\Crefname{originallemma}{Original lemma}{Original lemmas}
\makeatletter
\newlabel{lemma-norm-down-tame-symbol}{{0EAT}{source}{}{sourcepointer-5}{}}
\newlabel{lemma-norm-down-tame-symbol@cref}{{[originallemma][0][]0EAT}{[1][0][]source}{}{}{}}
\makeatother
\crefname{originallemma}{original lemma}{original lemmas}
\Crefname{originallemma}{Original lemma}{Original lemmas}
\makeatletter
\newlabel{lemma-perpare-key}{{0EAV}{source}{}{sourcepointer-6}{}}
\newlabel{lemma-perpare-key@cref}{{[originallemma][0][]0EAV}{[1][0][]source}{}{}{}}
\makeatother
\crefname{originallemma}{original lemma}{original lemmas}
\Crefname{originallemma}{Original lemma}{Original lemmas}
\makeatletter
\newlabel{lemma-powers-period-length-zero}{{0EAB}{source}{}{sourcepointer-7}{}}
\newlabel{lemma-powers-period-length-zero@cref}{{[originallemma][0][]0EAB}{[1][0][]source}{}{}{}}
\makeatother
\crefname{originallemma}{original lemma}{original lemmas}
\Crefname{originallemma}{Original lemma}{Original lemmas}
\makeatletter
\newlabel{lemma-prepare-tame-symbol}{{0EAG}{source}{}{sourcepointer-8}{}}
\newlabel{lemma-prepare-tame-symbol@cref}{{[originallemma][0][]0EAG}{[1][0][]source}{}{}{}}
\makeatother
\title{Vanishing of the sum of the orders of the tame symbols over the height one primes of a two-dimensional Noetherian local ring}
\author{The Stacks Project Authors}
\date{stacks-project:chow.tex:local-snapshot-20261012; GNU FDL 1.2 or later}
\begin{document}
\maketitle
\noindent\textit{Editorial note.} All delivered mathematical text is the author's own, copied from the pinned Stacks Project chapter file chow.tex; no statement or proof step was rewritten, reordered or removed.\par\medskip
\noindent\textit{Editorial note.} Preamble provenance: the pinned chapter file opens with \textbackslash{}input\{preamble\} and carries no preamble text of its own, so the project's single Stacks-wide preamble staged at sources/stacks-chow/preamble.tex is excerpted verbatim as two spans, lines 14-18 (Xy-pic setup) and lines 154-217 (the native theorem declarations and mathematical macros). Nothing was added, removed or redefined; the wrapper class and the amsmath/amssymb/amsthm/graphicx base packages are supplied by the assembler.\par\medskip
\noindent\textit{Editorial note.} Same-chapter pointers are located at their \textbackslash{}label lines inside the pinned chapter file; cross-chapter pointers are located in the project's own tag list staged in the snapshot (sources/stacks-chow/tags.txt), at the verbatim lines carrying their canonical tags.\par\medskip
\noindent\textit{Allowed standard prerequisites (editorial).}\begin{itemize}
\item Preparation of the tame symbol along each height one prime. Same chapter chow.tex, labels lemma-prepare-tame-symbol (proof line 1153) and lemma-perpare-key (proof line 1154), the latter applied with t = ab. Produces, for each height one prime q\_i, a finite local extension A inside B with B/A killed by a power of ab and B\_\{q\_i\} = B\_i; passing to B is the first reduction and is carried out in the counted span.
\item Norm compatibility of the tame symbol along a finite extension. Same chapter chow.tex, label lemma-norm-down-tame-symbol (proof line 1166); Algebra, label algebra-lemma-finite-extension-dim-1 (proof line 1167). Expresses the order of the tame symbol at q\_i as the sum of the orders at the primes of B over q\_i, so that the replacement of A by B preserves the total sum.
\item Comparison of periodic lengths and the multiplication formula. Same chapter chow.tex, labels lemma-compare-periodic-lengths (proof line 1195), lemma-multiply-period-length (proof line 1218), lemma-powers-period-length-zero (proof line 1225) and lemma-length-multiplication (proof lines 1227, 1241, 1257). Turns the vanishing of the sum of the e\_A(M\_i, a, b) into the statement about orders, and supplies the term by term identities used in the computation; the computation itself is inside the counted span.
\end{itemize}\medskip
\ifdefined\hypertarget\hypertarget{sourcepointer-1}{}\fi
\noindent\textit{Source reference.} Original lemma 02MJ: sources/stacks-chow/tags.txt, source lines 3446--3446. Navigation pointer for \textbackslash{}ref\{algebra-lemma-finite-extension-dim-1\} used inside the counted span (line 1167); the argument that consumes it is retained in the span.\par\medskip
\ifdefined\hypertarget\hypertarget{sourcepointer-2}{}\fi
\noindent\textit{Source reference.} Original lemma 0EA9: sources/stacks-chow/pinned-original-source.tex, source lines 283--283. Navigation pointer for \textbackslash{}ref\{lemma-compare-periodic-lengths\} used inside the counted span (line 1195); the argument that consumes it is retained in the span.\par\medskip
\ifdefined\hypertarget\hypertarget{sourcepointer-3}{}\fi
\noindent\textit{Source reference.} Original lemma 02QF: sources/stacks-chow/pinned-original-source.tex, source lines 310--310. Navigation pointer for \textbackslash{}ref\{lemma-length-multiplication\} used inside the counted span (line 1227); the argument that consumes it is retained in the span.\par\medskip
\ifdefined\hypertarget\hypertarget{sourcepointer-4}{}\fi
\noindent\textit{Source reference.} Original lemma 0EAC: sources/stacks-chow/pinned-original-source.tex, source lines 464--464. Navigation pointer for \textbackslash{}ref\{lemma-multiply-period-length\} used inside the counted span (line 1218); the argument that consumes it is retained in the span.\par\medskip
\ifdefined\hypertarget\hypertarget{sourcepointer-5}{}\fi
\noindent\textit{Source reference.} Original lemma 0EAT: sources/stacks-chow/pinned-original-source.tex, source lines 952--952. Navigation pointer for \textbackslash{}ref\{lemma-norm-down-tame-symbol\} used inside the counted span (line 1166); the argument that consumes it is retained in the span.\par\medskip
\ifdefined\hypertarget\hypertarget{sourcepointer-6}{}\fi
\noindent\textit{Source reference.} Original lemma 0EAV: sources/stacks-chow/pinned-original-source.tex, source lines 1105--1105. Navigation pointer for \textbackslash{}ref\{lemma-perpare-key\} used inside the counted span (line 1154); the argument that consumes it is retained in the span.\par\medskip
\ifdefined\hypertarget\hypertarget{sourcepointer-7}{}\fi
\noindent\textit{Source reference.} Original lemma 0EAB: sources/stacks-chow/pinned-original-source.tex, source lines 404--404. Navigation pointer for \textbackslash{}ref\{lemma-powers-period-length-zero\} used inside the counted span (line 1225); the argument that consumes it is retained in the span.\par\medskip
\ifdefined\hypertarget\hypertarget{sourcepointer-8}{}\fi
\noindent\textit{Source reference.} Original lemma 0EAG: sources/stacks-chow/pinned-original-source.tex, source lines 642--642. Navigation pointer for \textbackslash{}ref\{lemma-prepare-tame-symbol\} used inside the counted span (line 1153); the argument that consumes it is retained in the span.\par\medskip
\begin{lemma}
\label{lemma-key-nonzerodivisors}
Let $(A, \mathfrak m)$ be a $2$-dimensional Noetherian local ring.
Let $a, b \in A$ be nonzerodivisors.
Then we have
$$
\sum
\text{ord}_{A/\mathfrak q}(\partial_{A_{\mathfrak q}}(a, b))
=
0
$$
where the sum is over the height $1$ primes $\mathfrak q$ of $A$.
\end{lemma}

\begin{proof}
If $\mathfrak q$ is a height $1$ prime of $A$ such that $a, b$
map to a unit of $A_\mathfrak q$, then $\partial_{A_\mathfrak q}(a, b) = 1$.
Thus the sum is finite. In fact, if
$V(ab) = \{\mathfrak m, \mathfrak q_1, \ldots, \mathfrak q_r\}$
then the sum is over $i = 1, \ldots, r$.
For each $i$ we pick an extension $A_{\mathfrak q_i} \subset B_i$
as in Lemma \ref{lemma-prepare-tame-symbol} for $a, b$.
By Lemma \ref{lemma-perpare-key} with $t = ab$ and the given list of primes
we may assume we have a finite local extension $A \subset B$
with $B/A$ annihilated by a power of $ab$ and such that
for each $i$ the $B_{\mathfrak q_i} \cong B_i$.
Observe that if $\mathfrak q_{i, j}$ are the primes of $B$
lying over $\mathfrak q_i$ then we have
$$
\text{ord}_{A/\mathfrak q_i}(\partial_{A_{\mathfrak q_i}}(a, b))
=
\sum\nolimits_j
\text{ord}_{B/\mathfrak q_{i, j}}(\partial_{B_{\mathfrak q_{i, j}}}(a, b))
$$
by Lemma \ref{lemma-norm-down-tame-symbol} and
Algebra, Lemma \ref{algebra-lemma-finite-extension-dim-1}.
Thus we may replace $A$ by $B$ and
reduce to the case discussed in the next paragraph.

\medskip\noindent
Assume for each $i$ there is a nonzerodivisor
$\pi_i \in A_{\mathfrak q_i}$ and units $u_i, v_i \in A_{\mathfrak q_i}$
such that for some integers $e_i, f_i \geq 0$ we have
$$
a = u_i \pi_i^{e_i},\quad b = v_i \pi_i^{f_i}
$$
in $A_{\mathfrak q_i}$. Setting
$m_i = \text{length}_{A_{\mathfrak q_i}}(A_{\mathfrak q_i}/\pi_i)$
we have
$\partial_{A_{\mathfrak q_i}}(a, b) =
((-1)^{e_if_i}u_i^{f_i}v_i^{-e_i})^{m_i}$ by definition.
Since $a, b$ are nonzerodivisors the
$(2, 1)$-periodic complex $(A/(ab), a, b)$ has vanishing cohomology.
Denote $M_i$ the image of $A/(ab)$ in $A_{\mathfrak q_i}/(ab)$.
Then we have a map
$$
A/(ab) \longrightarrow \bigoplus M_i
$$
whose kernel and cokernel are supported in $\{\mathfrak m\}$
and hence have finite length. Thus we see that
$$
\sum e_A(M_i, a, b) = 0
$$
by Lemma \ref{lemma-compare-periodic-lengths}. Hence it suffices to show
$e_A(M_i, a, b) =
- \text{ord}_{A/\mathfrak q_i}(\partial_{A_{\mathfrak q_i}}(a, b))$.

\medskip\noindent
Let us prove this first, in case $\pi_i, u_i, v_i$ are the images
of elements $\pi_i, u_i, v_i \in A$ (using the same symbols
should not cause any confusion). In this case we get
\begin{align*}
e_A(M_i, a, b) & =
e_A(M_i, u_i\pi_i^{e_i}, v_i\pi_i^{f_i}) \\
& =
e_A(M_i, \pi_i^{e_i}, \pi_i^{f_i}) -
e_A(\pi_i^{e_i}M_i, 0, u_i) +
e_A(\pi_i^{f_i}M_i, 0, v_i) \\
& =
0 -
f_im_i\text{ord}_{A/\mathfrak q_i}(u_i) +
e_im_i\text{ord}_{A/\mathfrak q_i}(v_i) \\
& =
-m_i\text{ord}_{A/\mathfrak q_i}(u_i^{f_i}v_i^{-e_i}) =
-\text{ord}_{A/\mathfrak q_i}(\partial_{A_{\mathfrak q_i}}(a, b))
\end{align*}
The second equality holds by Lemma \ref{lemma-multiply-period-length}.
Observe that
$M_i \subset (M_i)_{\mathfrak q_i} = A_{\mathfrak q_i}/(\pi_i^{e_i + f_i})$
and
$(\pi_i^{e_i}M_i)_{\mathfrak q_i} \cong A_{\mathfrak q_i}/\pi_i^{f_i}$ and
$(\pi_i^{f_i}M_i)_{\mathfrak q_i} \cong A_{\mathfrak q_i}/\pi_i^{e_i}$.
The $0$ in the third equality comes from
Lemma \ref{lemma-powers-period-length-zero}
and the other two terms come from
Lemma \ref{lemma-length-multiplication}.
The last two equalities follow from multiplicativity of
the order function and from the definition of our tame symbol.

\medskip\noindent
In general, we may first choose $c \in A$, $c \not \in \mathfrak q_i$
such that $c\pi_i \in A$. After replacing $\pi_i$ by $c\pi_i$
and $u_i$ by $c^{-e_i}u_i$ and $v_i$ by $c^{-f_i}v_i$
we may and do assume $\pi_i$ is in $A$.
Next, choose an $c \in A$, $c \not \in \mathfrak q_i$
with $cu_i, cv_i \in A$. Then we observe that
$$
e_A(M_i, ca, cb) = e_A(M_i, a, b) - e_A(aM_i, 0, c) + e_A(bM_i, 0, c)
$$
by Lemma \ref{lemma-length-multiplication}.
On the other hand, we have
$$
\partial_{A_{\mathfrak q_i}}(ca, cb) =
c^{m_i(f_i - e_i)}\partial_{A_{\mathfrak q_i}}(a, b)
$$
in $\kappa(\mathfrak q_i)^*$ because $c$ is a unit in $A_{\mathfrak q_i}$.
The arguments in the previous paragraph show that
$e_A(M_i, ca, cb) = -
\text{ord}_{A/\mathfrak q_i}(\partial_{A_{\mathfrak q_i}}(ca, cb))$.
Thus it suffices to prove
$$
e_A(aM_i, 0, c) = \text{ord}_{A/\mathfrak q_i}(c^{m_if_i})
\quad\text{and}\quad
e_A(bM_i, 0, c) = \text{ord}_{A/\mathfrak q_i}(c^{m_ie_i})
$$
and this follows from Lemma \ref{lemma-length-multiplication}
by the description (see above)
of what happens when we localize at $\mathfrak q_i$.
\end{proof}

\end{document}
